\section{MARCO TEÓRICO}

\subsection{Web Semántica}

La Web Semántica es una extensión de la Web tradicional orientada a que los datos no
solo sean publicados para consumo humano, sino también descritos de manera formal para
ser interpretados por agentes de software \parencite{berners-lee2001}. Su objetivo no
es reemplazar la Web documental, sino complementarla con una capa semántica que permita
representar entidades, relaciones, restricciones y reglas de inferencia sobre los
recursos disponibles.

Desde el punto de vista arquitectónico, la Web Semántica se apoya en una pila de
estándares del W3C. En la base se encuentran los URI como identificadores globales de
recursos; sobre ellos, RDF como modelo de representación; posteriormente, RDFS y OWL
como lenguajes de descripción y formalización de vocabularios; y finalmente SPARQL
como lenguaje de consulta sobre grafos RDF. Gracias a esta pila tecnológica es posible
publicar conocimiento de manera interoperable y vincularlo con otras fuentes dentro del
ecosistema de \textit{Linked Data}.

\subsection{RDF como modelo de representación}

El \textit{Resource Description Framework} (RDF) es el modelo de datos básico de la Web
Semántica \parencite{w3c-rdf2014}. RDF expresa la información mediante triples de la
forma sujeto--predicado--objeto. El sujeto representa el recurso descrito, el
predicado expresa una propiedad o relación, y el objeto representa el valor de esa
propiedad, ya sea otro recurso o un literal.

Un conjunto de triples conforma un grafo RDF. Esta estructura es flexible, extensible y
adecuada para dominios complejos, ya que permite agregar nuevos hechos sin rediseñar el
modelo completo. En aplicaciones como la desarrollada en este trabajo, RDF es
particularmente útil porque permite integrar en un mismo grafo información sobre
propiedades, amenidades, perfiles de huésped, zonas geográficas y recursos enlazados a
fuentes externas como DBpedia.

\subsection{RDFS y anotaciones semánticas}

RDF Schema extiende RDF con construcciones básicas para describir tipos y jerarquías
\parencite{w3c-rdfs2014}. Entre sus elementos más importantes se encuentran
\texttt{rdfs:Class}, \texttt{rdfs:subClassOf}, \texttt{rdfs:domain} y
\texttt{rdfs:range}. Estas construcciones permiten definir clases, organizar
subsunciones entre ellas y restringir el uso de las propiedades.

Además, RDFS introduce anotaciones como \texttt{rdfs:label} y \texttt{rdfs:comment},
que permiten asociar nombres legibles y descripciones textuales a recursos. Estas
anotaciones son especialmente relevantes para el presente proyecto, ya que constituyen
la base del soporte multilingüe implementado en la ontología mediante etiquetas en
español, inglés y francés.

\subsection{Ontologías OWL}

Una ontología es una especificación formal de un dominio de conocimiento, donde se
declaran conceptos, relaciones, individuos y restricciones lógicas. OWL es el estándar
del W3C para la definición de ontologías sobre RDF \parencite{w3c-owl2012}. Frente a la
expresividad limitada de RDFS, OWL incorpora constructores más ricos para representar
disjunciones, equivalencias, restricciones existenciales, cardinalidades y axiomas
lógicos con semántica formal basada en lógica descriptiva.

En OWL, una ontología puede incluir:

\begin{itemize}
  \item \textbf{Clases}, que agrupan individuos con características comunes.
  \item \textbf{Propiedades de objeto}, que vinculan individuos con otros individuos.
  \item \textbf{Propiedades de dato}, que vinculan individuos con literales.
  \item \textbf{Individuos}, que representan instancias concretas del dominio.
  \item \textbf{Axiomas}, que describen restricciones, equivalencias, disjunciones o
  condiciones de pertenencia.
\end{itemize}

La ontología desarrollada en esta práctica fue modelada en OWL 2 DL, ya que este
perfil ofrece un equilibrio adecuado entre expresividad y decidibilidad. Gracias a ello,
el modelo puede ser verificado por un razonador automático y al mismo tiempo mantener
capacidad suficiente para representar el dominio de alojamientos temporales.

\subsection{SPARQL como lenguaje de consulta}

SPARQL es el lenguaje estándar para consultar grafos RDF \parencite{w3c-sparql2013}.
Una consulta SPARQL declara patrones de triples que deben satisfacerse en el grafo de
conocimiento y puede complementarse con filtros, uniones, agregaciones, ordenamiento y
agrupamiento de resultados.

En el contexto del presente trabajo, SPARQL se utiliza en dos niveles. Primero, como
mecanismo principal de recuperación sobre el dataset local cargado en Fuseki. Segundo,
como medio para demostrar interoperabilidad con DBpedia mediante consultas federadas con
la cláusula \texttt{SERVICE}. Aunque la aplicación final se apoya sobre un dataset local
enriquecido para mejorar estabilidad y rendimiento, el uso de \texttt{SERVICE} sigue
siendo relevante como evidencia de integración con una base de conocimiento enlazada.

\subsection{Lógica de Descripción y razonamiento}

La base formal de OWL 2 DL se encuentra en la Lógica de Descripción. Este marco permite
definir conceptos de manera precisa y habilita inferencias automáticas sobre la
ontología \parencite{baader2003}. A partir de un conjunto de axiomas, un razonador
puede verificar consistencia, clasificar jerarquías de clases y determinar la
pertenencia de individuos a conceptos derivados.

El razonamiento automático es importante en ontologías aplicadas porque permite detectar
errores conceptuales tempranos. Por ejemplo, si dos clases son declaradas disjuntas, un
individuo no debería pertenecer simultáneamente a ambas; si esto ocurre, el razonador lo
identifica como inconsistencia. En este proyecto se utilizó HermiT para verificar que la
ontología modelada no incurre en contradicciones lógicas \parencite{glimm2014}.

\subsection{Ingeniería de ontologías y preguntas de competencia}

La ingeniería de ontologías es el proceso metodológico mediante el cual se construyen y
mantienen modelos formales de conocimiento. Una práctica central en este proceso es la
definición de \textit{preguntas de competencia}: consultas representativas del dominio
que la ontología debe ser capaz de responder. Estas preguntas permiten validar la
cobertura conceptual del modelo y orientar su diseño desde necesidades reales de uso.

En el dominio de alojamiento temporal, las preguntas de competencia suelen combinar
ciudad, rango de precio, tipo de propiedad, amenidades, perfil del viajero y propósito
de la estadía. Por ello, la ontología no debe limitarse a almacenar atributos, sino que
debe representar relaciones semánticas que permitan responder consultas compuestas de
forma coherente.

\subsection{Protégé y HermiT como entorno de modelado}

Protégé es el entorno de edición ontológica más utilizado en proyectos de Web
Semántica \parencite{musen2015}. Permite modelar clases, propiedades, restricciones e
individuos, y además ofrece integración con razonadores como HermiT. En este trabajo,
Protégé fue utilizado para el diseño inicial de la ontología y HermiT para la
verificación de consistencia y revisión de la jerarquía lógica.

\subsection{DBpedia como base de conocimiento enlazada}

DBpedia es una de las bases de conocimiento enlazadas más representativas del ecosistema
de \textit{Linked Data}. Extrae información estructurada de Wikipedia y la publica en
RDF, con millones de entidades y etiquetas multilingües \parencite{lehmann2015}. Su
utilidad en este proyecto radica en que provee referencias geográficas reales para el
territorio boliviano y recursos equivalentes para conceptos del dominio como tipos de
alojamiento o amenidades.

La integración con DBpedia puede realizarse de varias maneras: mediante consultas
federadas en tiempo de ejecución, mediante descarga previa de subconjuntos RDF o a
través de enlaces semánticos entre recursos locales y remotos. En la solución final se
priorizó este último enfoque, incorporando \texttt{owl:sameAs} y
\texttt{owl:equivalentClass} en la ontología local, lo que hace más estable la
demostración sobre Fuseki y, al mismo tiempo, conserva la trazabilidad hacia los
recursos de DBpedia utilizados.

\subsection{Representación multilingüe en ontologías}

Una de las ventajas más importantes de RDF/OWL es que la multilingüidad puede
representarse directamente mediante literales etiquetados por idioma. La forma más
simple y estándar consiste en utilizar \texttt{rdfs:label} y \texttt{rdfs:comment}
con etiquetas como \texttt{@es}, \texttt{@en} o \texttt{@fr} \parencite{w3c-rdf2014}.

Este enfoque tiene varias ventajas técnicas: evita duplicar la ontología, mantiene una
única identidad semántica para cada recurso y permite seleccionar el idioma requerido en
las consultas SPARQL mediante la función \texttt{LANG()} o estrategias de fallback.
Precisamente esta capacidad fue aprovechada en el proyecto para servir resultados en
español, inglés y francés desde un único dataset OWL cargado en Fuseki.
